% 3DV 2025 Paper Template
\documentclass[10pt,twocolumn,letterpaper]{article}

% \usepackage{cvpr}              % camera-ready
\usepackage[review]{cvpr}      % review version
% \usepackage[pagenumbers]{cvpr}

\input{preamble}

\definecolor{cvprblue}{rgb}{0.21,0.49,0.74}
\usepackage[pagebackref,breaklinks,colorlinks,citecolor=cvprblue]{hyperref}

\def\paperID{*****}
\def\confName{3DV\xspace}
\def\confYear{2025\xspace}

\title{Benchmarking Spatial and State Reasoning Skills of\\
Vision-Language Models for Robotics}

\author{Alexander Padula$^1$ \quad Anna Kravchenko$^1$ \quad Christian Deubel$^1$ \quad Defne Soganci$^1$\\
Jovana Videnovi\'{c}$^2$ \quad Jelena Trisovic$^2$ \quad Rene Zurbrugg$^2$\\[4pt]
$^1$ETH Zurich, D-INFK \quad $^2$Computer Vision and Geometry Lab
}

\begin{document}
\maketitle

% Abstract should answer:
% - What problem are we tackling (VLM spatial/state reasoning for robotics)?
% - What benchmark/approach did we build to test it?
% - What are the headline findings (zero-shot gap, CoT hurts, LoRA gain decomposition)?
% LENGTH TARGET: ~150-200 words, single paragraph (CVPR norm)
\begin{abstract}
Vision-Language Models (VLMs) have been proposed as action-prediction models and
robot supervisors, roles that both demand reliable judgment of task progress and failure
detection from camera input.
We introduce a new benchmark targeting failure-mode detection and multi-view answer consistency. It comprises 6224 questions across 47 hand-labeled scenes
drawn from Robo2VLM-1 (DROID).
We evaluate seven VLMs from three families (Gemma3, Qwen3, InternVL3) zero-shot,
with a chain-of-thought (CoT) ablation and a LoRA fine-tuning case study on Qwen3-4B.
Zero-shot accuracy sits only marginally above random (high-40s\% vs.\ 33\%) and plateaus
at 30B+ scale; CoT prompting consistently hurts overall accuracy; and LoRA fine-tuning
yields a +10.9pp headline gain, but roughly 80\% of the gain can be attributed to label-position bias and learning to detect out-of-distribution scenes rather than improved temporal or task reasoning.
We find that current small- and mid-sized open-weight VLMs perform remarkably poorly as robot supervisors, and
fine-tuning does not significantly improve their performance.
\end{abstract}


% Introduction should answer:
% - Why are VLMs relevant as zero-shot "robot supervisors" (scene understanding + feedback
%   without task-specific training)?
% - Can a VLM judge task progress, current action phase, and success/failure from a single
%   image + stated goal?
% - Do answers extracted from different camera angles of the same moment agree with each
%   other?
% - What is the contribution summary (benchmark, models evaluated, ablations, LoRA case study)?
\section{Introduction}
\label{sec:intro}

Recent advances in Vision-Language-Models~(VLM) brought LLM-type scaling and reasoning
into the domain of robotics. As VLMs are used in both next-action prediction
(Vision-Language-Action models) and higher-level task planning and supervision, they
require state understanding of the robot and the objects it interacts with, given the
context of the task. To bridge the gap to real-world deployment, the models have to be 
self-sufficient, thus generalizing to unforeseen
environments, tasks, and edge cases, relieving the user from needing to finetune the
model excessively. Thus it is necessary for these models to reliably track progress during a
task, identify which sub-task the robot is currently attempting, and understand
success/failure of robotic actions. Furthermore it has to be time and resource efficient
based on the task and an overall pipeline understanding: e.g.\ for a household robot it
might suffice to empty the dishwasher under 20 minutes, while an industrial production
pipeline might have only several seconds to determine success or failure of a task. We
investigate the zero-shot capabilities of different model families and sizes for failure
mode understanding and multi-view consistency, to determine under which task-specific and
scene conditions the models' reasoning collapses.

%TODO
We find that current small- and mid-sized VLMs fail to deliver this level of reliability:
zero-shot accuracy stays close to random across model families, and explicitly prompting
models to reason step by step before answering makes performance worse, not better. While
such models exhibit robust performance across object detection, positional understanding, they struggle with object states and complex cross-object relationships.  robot supervision poses a
narrower and harder questions: Robots require supervision beyond raw
perception: detecting failures, tracking progress, and reasoning about the current action
phase are prerequisites for any higher-level planning or human-in-the-loop monitoring
built on top of a perception stack. The open question we address is at which task- and scene-complexity, models fail. 

We study this gap along two concrete capability questions. First, failure-mode and state
reasoning: given multiple views of the scene, can a VLM judge task progress,
identify the current action phase, and determine whether action-phase execution succeeded or failed?
Second, multi-view consistency: do answers extracted from different camera angles of the
 moment agree with each other? Inconsistency across views would
undermine true scene understanding and entails intensive calibration phase.

To answer these questions, we build a new benchmark on
Robo2VLM-1~\cite{chen2025robo2vlm} (DROID~\cite{khazatsky2024droid} and OXE
trajectories), covering 47 scenes and 6224 questions across four question types: four
failure-mode sub-tasks and one multi-view consistency task. We evaluate seven VLMs across
three families zero-shot, ablate chain-of-thought (CoT) prompting, analyze accuracy as a
function of task- and scene-difficulty, and run a LoRA fine-tuning case study on Qwen3-4B with a
decomposition of where the performance gain actually come from.
\begin{itemize}
  \item New benchmark: 6224 questions / 47 scenes on Robo2VLM-1 (DROID), covering
        failure-mode detection (4 sub-tasks) and multi-view consistency.
  \item Self-designed question types (action phase, phase success, progress, task success,
        multi-view consistency), including adversarial ``Cannot be determined'' (CnBD) variant.
  \item Zero-shot evaluation of 7 VLMs across 3 families (Gemma3, Qwen3, InternVL3) with
        easy/hard scene-difficulty breakdown and CoT ablation.
  \item LoRA fine-tuning case study (Qwen3-4B): +10.9pp overall, with detailed gain
        decomposition revealing two coarse heuristics rather than improved phase reasoning.
\end{itemize}

% --- Figure: method overview ---
\begin{figure}[t]
  \centering
  \includegraphics[width=\linewidth]{../poster_3_D/method_overview.pdf}
  \caption{Benchmark pipeline: trajectories from Robo2VLM-1 (DROID/OXE) are turned into
  multiple-choice failure-mode and multi-view consistency questions, which are posed to
  seven zero-shot VLMs, a CoT variant, and a LoRA fine-tuned model, then scored and
  decomposed by question type and answer class.}
  \label{fig:method_overview}
\end{figure}


% Related Work should answer:
% - What existing datasets generate VQA from robot trajectories (Robo2VLM-1 and similar)?
% - How have VLMs been used as evaluators/supervisors/reward models in prior work?
% - What known VLM weaknesses exist in spatial/3D reasoning and chain-of-thought for visual
%   tasks?
\section{Related Work}
\label{sec:related}

\textbf{VQA from robot trajectories.} Robo2VLM-1~\cite{chen2025robo2vlm} automatically
generates visual-question-answering pairs from real robot trajectories using
DROID~\cite{khazatsky2024droid} and Open X-Embodiment (OXE) data, deriving ground-truth
answers from sensor and trajectory metadata rather than human annotation; this work
builds directly on it. They leverage the multiple synchronizes views per trajectory in DROID~\cite{khazatsky2024droid} (top-left, top-right, wrist), which is what
enables our multi-view consistency study. Prior robotics VQA benchmarks scaled poorly and suffered from scene homogeneity. 

\textbf{VLMs as evaluators and robot supervisors.} VLMs have been proposed for action-prediction, as well as supervision, where supervision can be viewed as a mandatory subtask of the former. As standalone supervision models, they are used as judges for production outcomes and high level planning to execute the gating logic for the next action phases. 
detectors, reward models, and failure classifiers for robotic manipulation
\TODO{cite RT-2 / VoxPoser / success-detector papers}, and a separate line of work studies
VLMs as general-purpose judges, showing that judgment quality depends heavily on prompt
format, answer ordering, and question complexity \TODO{cite VLM-as-judge literature}. Our
work differs in two ways: we evaluate both failure-mode detection and multi-view
consistency on the same real-robot data, and we decompose fine-tuning gains by question
type and answer class rather than reporting a single headline accuracy number.

\textbf{Known VLM weaknesses in spatial and temporal reasoning.} Survey evidence suggests
VLMs trained on web-scale data overfit to object-identity and object-position queries,
while state-change and contact reasoning are comparatively underrepresented in their
training distributions~\cite{lin2025vlmsurvey}. Evidence on chain-of-thought prompting for
visual tasks is mixed: it helps on math- and logic-heavy benchmarks but can hurt on
perceptual tasks by introducing hallucinated reasoning chains
\TODO{cite CoT-hurts-perception literature}; our CoT ablation
(\cref{sec:results}) adds a robotics data point to this debate. The model families we
evaluate are Gemma3~\cite{gemma2025gemma3}, Qwen3~\cite{qwen2025qwen3}, and
InternVL3~\cite{chen2025internvl3}.


% Dataset & Tasks should answer:
% - Where does Robo2VLM-1 come from (DROID/OXE) and how are questions/images generated?
% - What are the 4 action-phase sub-tasks, and what is the CnBD adversarial variant?
% - How does the multi-view benchmark define cross-camera consistency?
% - What are the dataset statistics (scenes, questions per task, train/val/test splits)?
% - What baselines is model performance compared against?
\section{Dataset \& Tasks}
\label{sec:dataset}

\subsection{Source Data}
To leverage Robo2VLM-1~\cite{chen2025robo2vlm} keyframe selection of the scene diverse DROID~\cite{khazatsky2024droid} dataset.
\cite{chen2025robo2vlm} use a novel framework to identify keyframes within each
trajectory and pair them with the goal text, ordered action-phase labels, and
sensor-derived success/failure signals captured from three synchronized camera views
(top-left, top-right, wrist). From this pool we select 47 scenes from the DROID subset,
covering diverse manipulation tasks such as placing, screwing, folding, and stacking. We label and sort the action phases manually to ensure (1) high quality and diverse scene selection (2) scene and task understanding is humanly possible (3) securing groundtruth annotation on which view holds necesssary visual signal (4) temporal consistency.
As a result we obtain an ordered list of keyframes, their action-phase labels of what the robot is doing, as well as annotation of which view enables comprehension and correct answering. 

We define four multiple-choice failure-mode sub-tasks, all evaluated on real DROID trajectory keyframes.

\textbf{Action Phase Identification} (\texttt{action\_phase\_id}): given a single moment, the goal text, and a claimed current action phase, the model selects which phase (A--D) is currently depicted, or answers ``Cannot be determined'' (E). Random baseline: 33.3\%. 

\textbf{Phase Success} (\texttt{phase\_success}): given a single moment, the goal text, and
a target phase, the model judges whether that phase is Before, Is currently performing,
After, or Cannot be determined. Random baseline: 25.0\%. 

\textbf{Progress Detection} (\texttt{progress}): given two distinct moments and the
goal text, the model judges whether the task progressed forward, backward, or
Cannot be determined. Random baseline: 33.3\%.

\textbf{Task Success} (\texttt{task\_success}): given a single image and the goal text,
the model answers Yes (complete), No (failed or incomplete), or Cannot be determined.
Random baseline: 33.3\%.

For \texttt{action\_phase\_id}, \texttt{task\_success}, and \texttt{progress} we
additionally inject an adversarial \emph{Cannot be determined} (CnBD) variant: when
\texttt{special\_image == "random\_scene"}, the input image is replaced with an unrelated
scene, making ``Cannot be determined'' the only correct answer. This variant tests
whether a model detects a domain mismatch between image and goal text, rather than
guessing a plausible-looking phase.
Additionally we provide each question with context of a sorted list of action phase labels of the scene, to provide the model with potentially needed context. E.g. for action phase detection, the model might answer incorrectly, because it would attribute it sligthley differently formulated action phase.


% Figure: evaluation target examples
\begin{figure}[t]
  \centering
  \includegraphics[width=\linewidth]{../poster_3_D/phase_success_top.pdf}
  \caption{Example questions for the four failure-mode detection sub-tasks.
  \TODO{Write caption describing what each row shows.}}
  \label{fig:task_examples}
\end{figure}

\subsection{Failure-Mode Detection Sub-Tasks}
For failure detection, we prompt the model with all 3 views merged into a single image, as prodvided by \cite{chen2025robo2vlm}. CoT shows the models can identify the individual views in a merged image.
\subsection{Multi-View Consistency}
The same question is additionally posed from multiple camera views of the same scene, if a human labelled the view as "action-phase is comprehensible from single view":
top-left (TL), top-right (TR), and wrist (W). 
\begin{figure}[t]
  \centering
  \includegraphics[width=\linewidth]{../poster_3_D/mutlview_consistency_evaluation_targets_bottom.pdf}
  \caption{Multi-view consistency evaluation: same scene captured from top-left, top-right,
  and wrist cameras. \TODO{Write caption.}}
  \label{fig:multiview_examples}
\end{figure}

\subsection{Dataset Statistics}

\begin{table}[h]
  \centering
  \caption{Dataset statistics. The LoRA split uses a scene-level partition with no
  scene overlap between train, val, and test.}
  \label{tab:dataset_stats}
  \begin{tabular}{lrrr}
    \toprule
    & \textbf{Train} & \textbf{Val} & \textbf{Test} \\
    \midrule
    Scenes      & 20    & 7   & 20   \\
    Questions   & 2535  & 983 & 2706 \\
    \midrule
    \multicolumn{4}{l}{\small Total: 47 scenes, 6224 questions} \\
    \bottomrule
  \end{tabular}
\end{table}

Table~\ref{tab:dataset_stats} reports scene and question counts for the LoRA train/val/test
split. The full benchmark, used for zero-shot and CoT evaluation, totals 47 scenes and
6224 questions; the LoRA split is scene-disjoint, so no scene appears in more than one of
train, val, and test.
% TODO: Add per-sub-task question count breakdown if available.


% Method should answer:
% - Which models are evaluated, and how are they prompted (zero-shot multiple-choice)?
% - What does the chain-of-thought variant change about the prompting?
% - What is the LoRA fine-tuning setup (scene-level split, hyperparameters, eval pipeline)?
\section{Method}
\label{sec:method}

\subsection{Models}
We evaluate seven publicly available VLMs spanning three model families, chosen to cover
a range of parameter counts within and across families: Gemma3~\cite{gemma2025gemma3}
(E2B, E4B, 31B), Qwen3~\cite{qwen2025qwen3} (4B, 8B, 32B), and
InternVL3~\cite{chen2025internvl3} (14B). All models are run zero-shot on the ETH Zurich
student cluster (A100 GPUs) without any task-specific training.

\subsection{Zero-Shot Evaluation Setup}
The models are prompted with a single image (or two, for \texttt{progress}), the goal
text, and a multiple-choice question with answer options labeled A through E (the exact
count depends on the sub-task). We extract the first letter predicted in a single forward
pass as the model's answer, without few-shot examples or task-specific instructions. In
the chain-of-thought (CoT) variant, the model is instead asked to reason freely before
stating a final answer letter; we apply the same extraction procedure to the last letter
emitted in the response. \TODO{include the exact prompt template here, or move it to
supplementary.}

\subsection{LoRA Fine-Tuning}
We additionally fine-tune Qwen3-4B with LoRA to control for the possibility that low
zero-shot accuracy reflects missing task-specific grounding rather than a fundamental
reasoning ceiling. We use rank 16, alpha 32, and dropout 0.05 on the query, key, value,
and output projections of the language-model layers, trained on a scene-level split with
20 training scenes (2535 questions) and 7 validation scenes (983 questions); the held-out
test split (20 scenes, 2706 questions) shares no scene with training, preventing the
model from memorizing specific visual scenes. Training runs for 3 epochs with a learning
rate of $2\times10^{-4}$, batch size 1 with gradient accumulation 8 (effective batch size
8), and a maximum sequence length of 1024 tokens (seed 42). We compare against zero-shot
Qwen3-4B evaluated on the identical 2706-question test split, giving an apples-to-apples
BASE vs.\ +LoRA comparison.


% Results should answer:
% - How accurate are models at action-phase/failure-mode detection vs. random/majority
%   baselines?
% - How consistent are multi-view answers across camera angles?
% - Does chain-of-thought prompting help or hurt accuracy?
% - How does accuracy vary across scenes (difficulty spread)?
% - How much does LoRA fine-tuning improve accuracy, and what is the gain actually made of?
\section{Results}
\label{sec:results}

\subsection{Failure-Mode Detection: Zero-Shot}
Table~\ref{tab:main_results} reports zero-shot accuracy on the four failure-mode
sub-tasks and on multi-view consistency. All seven models score above random on
\texttt{action\_phase\_id} and \texttt{task\_success}, but sit much closer to random on
\texttt{phase\_success} and \texttt{progress} -- the two sub-tasks that require comparing
across phases or images rather than judging a single frame in isolation. Scale helps
within a family: Gemma3 improves monotonically from E2B to E4B to 31B on every sub-task
(e.g.\ 31.7\%~$\to$~41.4\%~$\to$~46.1\% on \texttt{action\_phase\_id}). \texttt{task\_success}
is an outlier: the Qwen3 family scores 70.9--72.6\%, nearly twice random, while Gemma3 and
InternVL3-14B plateau in the 45--59\% range, which we attribute to Qwen3's stronger
instruction following on a binary success/failure framing rather than better state
understanding. Even so, the best model on the hardest sub-task (Qwen3-32B, 48.0\% on
\texttt{action\_phase\_id}) is barely above the 33.3\% random baseline -- absolute
numbers tell a sobering story about current VLM state-reasoning capability. Consistent
with this, models rarely use ``Cannot be determined'' even where it is the correct
answer, defaulting instead to a confident Yes or No (\cref{sec:results-lora} quantifies
this for the CnBD variant under BASE Qwen3-4B).

\begin{table}[t]
  \centering
  \caption{Zero-shot accuracy (\%) on failure-mode detection sub-tasks and multi-view
  consistency (All\,3V). Best per column in \textbf{bold}.}
  \label{tab:main_results}
  \setlength{\tabcolsep}{4pt}
  \begin{tabular}{l|cccc|c}
    \toprule
    \textbf{Model} & \textbf{Act.Ph.} & \textbf{Ph.Succ.} & \textbf{Prog.} & \textbf{Task\,S.} & \textbf{All\,3V} \\
    \midrule
    Gemma3-E2B    & 31.7          & 24.6          & 35.0          & 37.4          & 27.9 \\
    Gemma3-E4B    & 41.4          & 30.0          & 38.4          & 45.7          & 45.2 \\
    Gemma3-31B    & 46.1          & \textbf{46.2} & \textbf{45.2} & 58.8          & 58.2 \\
    \midrule
    Qwen3-4B      & 44.8          & 35.6          & 39.1          & \textbf{72.6} & 51.0 \\
    Qwen3-8B      & 43.5          & 39.5          & 40.2          & 70.9          & \textbf{59.0} \\
    Qwen3-32B     & \textbf{48.0} & 42.6          & 40.9          & 72.2          & 49.2 \\
    \midrule
    InternVL3-14B & 43.8          & 43.8          & 41.7          & 56.2          & 56.1 \\
    \midrule
    Random        & 33.3          & 25.0          & 33.3          & 33.3          & 25.0 \\
    \bottomrule
    \multicolumn{6}{l}{\scriptsize Random: uniform (cols 1--4) / independent-binary (All\,3V).} \\
  \end{tabular}
\end{table}

\subsection{Multi-View Consistency}
Table~\ref{tab:multiview_results} reports the three consistency metrics. Qwen3-8B
(59.0\%), Gemma3-31B (58.2\%), and InternVL3-14B (56.1\%) are the most self-consistent
across views; Gemma3-E2B (27.9\%) is the weakest, barely above the 25.0\%
independent-binary random baseline. For every model, TL$\leftrightarrow$TR agreement is
higher than All-3-agree, indicating that the wrist view is the hardest to align with the
two top-down views. Breaking the (fully completed) Qwen3-32B run down by correctness
(Table~\ref{tab:multiview_correctness}) shows that models are most self-consistent when
uniformly correct (72.3\%) or uniformly wrong (66.4\%) across the three views, and least
consistent when only some views are correct (34.4\%) -- consistency tracks how
unambiguous a scene appears to the model, not whether its answer happens to be right.

\begin{table}[t]
  \centering
  \caption{Qwen3-32B all-3-agree rate (\%) by view-correctness group (1367 source groups).}
  \label{tab:multiview_correctness}
  \begin{tabular}{lc}
    \toprule
    \textbf{Group} & \textbf{All\,3V agree} \\
    \midrule
    All views correct  & 72.3 \\
    Some views correct & 34.4 \\
    No views correct   & 66.4 \\
    \bottomrule
  \end{tabular}
\end{table}

\begin{table}[t]
  \centering
  \caption{Multi-view consistency (\%) across three agreement metrics. Best per column in
  \textbf{bold}.}
  \label{tab:multiview_results}
  \setlength{\tabcolsep}{4pt}
  \begin{tabular}{l|ccc}
    \toprule
    \textbf{Model} & \textbf{TL$\leftrightarrow$TR} & \textbf{$\geq$1\,Top$\leftrightarrow$W} & \textbf{All\,3V} \\
    \midrule
    Gemma3-E2B    & 48.3          & 63.7          & 27.9 \\
    Gemma3-E4B    & 61.2          & 71.0          & 45.2 \\
    Gemma3-31B    & 75.9          & 78.9          & 58.2 \\
    \midrule
    Qwen3-4B      & 70.2          & 73.9          & 51.0 \\
    Qwen3-8B      & \textbf{74.8} & \textbf{81.6} & \textbf{59.0} \\
    Qwen3-32B     & 68.4          & 72.1          & 49.2 \\
    \midrule
    InternVL3-14B & 73.7          & 78.7          & 56.1 \\
    \midrule
    Random        & 50.0          & 75.0          & 25.0 \\
    \bottomrule
    \multicolumn{4}{l}{\scriptsize Random: independent per-view binary baseline.} \\
  \end{tabular}
\end{table}

\subsection{Chain-of-Thought Ablation}
Table~\ref{tab:cot_results} compares default vs.\ CoT prompting for the two Gemma3
models with complete CoT runs at time of writing (the Qwen3-4B CoT run had completed only
814 of 6224 questions, see \cref{sec:limitations}). CoT prompting consistently hurts:
Gemma3-E4B's \texttt{task\_success} accuracy drops from 45.7\% to 20.9\%, below the
33.3\% random baseline, and Gemma3-E2B's drops from 37.4\% to 15.5\%, with its overall
accuracy on the full 6224-question set falling to 28.8\%, below the 29.8\% random
baseline for that set. Two independent models now show CoT hurting accuracy, and one
drags overall accuracy below random -- this looks like a pattern rather than a single
model's anomaly. We hypothesize that models describe the visual content correctly but
draw an incorrect conclusion partway through the free-form reasoning chain, which then
anchors an incorrect final letter; confirming this would require an error-taxonomy pass
over the CoT transcripts, which we leave to future work (\cref{sec:limitations}).

\begin{table}[t]
  \centering
  \caption{Default vs.\ chain-of-thought (CoT) accuracy (\%) on \texttt{task\_success}.}
  \label{tab:cot_results}
  \begin{tabular}{lcc}
    \toprule
    \textbf{Model} & \textbf{Default} & \textbf{CoT} \\
    \midrule
    Gemma3-E2B & 37.4 & 15.5 \\
    Gemma3-E4B & 45.7 & 20.9 \\
    \midrule
    Random & 33.3 & 33.3 \\
    \bottomrule
  \end{tabular}
\end{table}

\subsection{Scene Difficulty Analysis}
Per-scene accuracy varies widely: the best zero-shot model spans roughly 31 percentage
points across the 47 scenes (Figure~\ref{fig:easy_hard}). Averaging across all seven
models, the consensus-hardest scene (\textit{align the rope along the middle of the
black object}, 35.1\% mean accuracy) sits 13.6pp below the consensus-easiest scene
(\textit{pick the orange ball and put it in the pot on the stove}, 48.7\%), and these
consensus rankings are far more stable than any single model's own ranking: the average
pairwise Spearman correlation of per-scene accuracy across the seven models is only
0.194, meaning models largely disagree on which scenes are hard. Qualitatively, easy
scenes feature visually distinct objects and clear state changes; hard scenes are
cluttered, have visually similar phases, or involve subtle state changes. This points to
scene visual properties -- clutter, contrast, object distinctness -- rather than task
semantics as the main driver of difficulty.

\begin{figure}[t]
  \centering
  \includegraphics[width=\linewidth]{../poster_3_D/easy_hard_scenes.pdf}
  \caption{Per-scene accuracy distribution (best zero-shot model). Easy scenes (left) feature
  distinct objects and clear state changes; hard scenes (right) are cluttered or visually
  ambiguous. \TODO{Write caption.}}
  \label{fig:easy_hard}
\end{figure}

\subsection{LoRA Fine-Tuning Case Study (Qwen3-4B)}
\label{sec:results-lora}
Table~\ref{tab:lora_results} compares Qwen3-4B BASE against +LoRA on the identical
2706-question, 20-scene held-out test split. Overall accuracy rises from 44.2\% to
55.1\% (+10.9pp, +295 correct answers) -- a substantial headline gain. Decomposing this
gain, however, shows that most of it is not improved phase or temporal reasoning but two
coarse heuristics.

The first heuristic (roughly 44\% of the gain) is improved detection of the adversarial
CnBD variant: on the 130 random-scene questions in \texttt{action\_phase\_id} and
\texttt{task\_success} (6.5\% of the test set), BASE accuracy is 0--15\% while +LoRA
reaches 78--88\%. The model has effectively learned the rule ``if the image scene does
not match the goal text, answer Cannot be determined.'' This is a real, generalizable
skill rather than memorization: splitting CnBD recall by the donor scene's origin gives
comparable recall whether the donor scene was seen in training (0.92 / 0.65 for
\texttt{action\_phase\_id} / \texttt{task\_success}), in validation (1.0 / 1.0), or
entirely unseen (0.82 / 0.82) -- recall on unseen donor scenes is not lower, so the rule
generalizes to novel scene pairings rather than memorizing specific images.

The second heuristic (roughly 55\% of the gain, 161 questions) appears on
\texttt{phase\_success} at MIDDLE phase positions (where the depicted image step equals
the target phase, i.e.\ neither first nor last): accuracy improves from 36.9\% to 49.9\%.
This gain, however, comes with a regression: +LoRA nearly abandons ``Is currently
performing it'' (IPI) as an answer, with IPI recall dropping from 35.8\% to 6.6\% and raw
IPI predictions falling from 383 to 43 across the full test set. Under a lenient scoring
rule that also accepts ``Yes'' for IPI, +LoRA is still worse than BASE (58.0\% vs.\
48.9\%), and of the 274 ground-truth IPI cases, +LoRA predicts ``No'' on 140 of them.
This is a net-positive headline number that is, in fact, a regression for the most
safety-critical use case: detecting that a robot is currently in the middle of performing
an action.

Notably, \texttt{progress} -- which also contains a CnBD variant -- does not pick up the
same OOD-detection cue despite comparable training exposure (110 vs.\ 130 examples),
which we attribute to \texttt{progress} swapping only one of its two side-by-side images
for the random scene, making the cue visually less salient than the full-image swap used
in \texttt{action\_phase\_id} and \texttt{task\_success}.

\begin{table}[t]
  \centering
  \caption{Qwen3-4B LoRA fine-tuning results on matched 20-scene / 2706-question
  held-out test split. $\Delta$ = LoRA minus BASE.}
  \label{tab:lora_results}
  \setlength{\tabcolsep}{4pt}
  \begin{tabular}{l|cccc|c}
    \toprule
    \textbf{Qwen3-4B} & \textbf{Act.Ph.} & \textbf{Ph.Succ.} & \textbf{Prog.} & \textbf{Task\,S.} & \textbf{Overall} \\
    \midrule
    BASE    & 45.1          & 37.9          & \textbf{38.2} & 73.5          & 44.2 \\
    +LoRA   & \textbf{61.0} & \textbf{51.8} & 37.5          & \textbf{81.1} & \textbf{55.1} \\
    \midrule
    $\Delta$ & +15.9        & +13.9         & $-$0.8        & +7.6          & +10.9 \\
    \bottomrule
  \end{tabular}
\end{table}

% TODO: Add gain-decomposition bar chart figure once generated.
% TODO: Add CnBD example (BASE wrong vs LoRA correct) figure.
% TODO: Add IPI-collapse example figure (ground truth IPI, BASE correct, LoRA wrong).

% Discussion should answer:
% - Why are absolute accuracies so close to random chance?
% - Does model scale help, and is it enough on its own?
% - Is the LoRA gain genuine phase/temporal reasoning or coarse heuristics (and is one of
%   those heuristics actually a trade-off)?
% - What drives multi-view consistency/inconsistency (confidence vs. correctness)?
\section{Discussion}
\label{sec:discussion}

\textbf{Why are absolute numbers so close to random?} Object \emph{state} -- deformation,
contact, orientation change -- is underrepresented in web-scale VLM training data
relative to object \emph{identity} and \emph{position}. Models can usually describe which
objects are present and roughly where, but not reliably what has happened to them or
which phase of an action is depicted. Compounding the problem, the models rarely use
``Cannot be determined'' as an answer even when it is correct, defaulting instead to a
confident Yes or No (\cref{sec:results-lora}: BASE Qwen3-4B scores 0--15\% on the CnBD
variant). The 70\%-plus \texttt{task\_success} accuracy of the Qwen3 family is a partial
exception: binary success/failure framing may be closer to the kind of outcome
classification VLMs see at web scale.

\textbf{Scale helps, but is not enough.} Within-family scaling consistently improves
accuracy (Gemma3 E2B~$\to$~E4B~$\to$~31B), but even 31B/32B-class models plateau in the
high-40s\% on phase-level reasoning. Cross-family, InternVL3-14B is competitive with
Qwen3-8B despite the parameter gap, suggesting architecture and training data matter at
least as much as raw parameter count. More parameters alone are unlikely to close the
gap; robot-domain training data is needed.

\textbf{LoRA gains require decomposition before deployment decisions.} The 11pp headline
improvement could be read as ``fine-tuning works for robot state reasoning,'' but the
decomposition in \cref{sec:results-lora} shows that roughly 80\% of the gain is two
coarse heuristics. The OOD-scene detector is genuinely useful and generalizable, but it is
domain-mismatch detection, not phase or temporal reasoning. The label-position bias on
\texttt{phase\_success} is net positive under this dataset's label distribution, but a
regression for the most safety-critical use case: detecting when a robot is currently in
the middle of an action. The lesson is to always decompose fine-tuning gains by question
type and answer class before concluding that a model is ``better'' at a target skill.

\textbf{Multi-view consistency tracks confidence, not correctness.} Models are most
self-consistent when uniformly right (72.3\%) or uniformly wrong (66.4\%) across views,
and least consistent when only some views are correct (34.4\%). We read this as
view-consistency reflecting how unambiguous a scene is to the model, not whether its
answer is correct. This connects to the scene-difficulty finding (\cref{sec:results}):
scenes that drive accuracy down also destabilize cross-view agreement, suggesting both
effects share the same root cause -- visual ambiguity, not task semantics. For multi-view
fusion, this implies that simple voting across views may not help on exactly the hardest
scenes, where inconsistency is highest.


% Conclusion should answer:
% - What are the three main findings to restate (marginal zero-shot gains, CoT hurts,
%   fine-tuning gains are mostly coarse heuristics)?
% - What is the single most important takeaway for the reader?
\section{Conclusion}
\label{sec:conclusion}

We benchmark VLMs on two robot-supervision capabilities and arrive at three findings.
First, zero-shot state reasoning is weak but non-trivial: accuracy stays only marginally
above random across model families, scale helps within a family, but even 30B+ models
plateau well short of reliable supervision. Second, chain-of-thought prompting hurts
rather than helps: two independent models drop below random on \texttt{task\_success}
with CoT, suggesting free-form reasoning chains introduce hallucinated conclusions on
perceptual tasks rather than improving them. Third, fine-tuning gains are real but mostly
coarse: LoRA lifts Qwen3-4B accuracy by +10.9pp, but roughly 80\% of that gain decomposes
into a generalizable out-of-distribution scene detector and a label-position bias that is
in fact a regression for the most critical failure-detection cases. Taken together, these
results suggest current VLMs are not yet ready for reliable robot-state supervision out
of the box, and that fine-tuning gains must be decomposed by question type and answer
class before they can be trusted for deployment.

\section{Limitations \& Future Work}
\label{sec:limitations}

\textbf{Limitations.} Manual annotation limits scale: 47 scenes is enough for a
controlled study but too small to draw strong statistical conclusions about
model-family differences; an automated annotation pipeline would let this scale up.
Train and test splits are drawn from the same Robo2VLM-1 distribution, so our results do
not test sim-to-real or cross-dataset generalization. Some OXE episodes cannot be traced
back to their raw source, limiting reproducibility for those questions. Two runs were
incomplete at the time of writing: the Qwen3-4B CoT run (814 of 6224 questions) and the
Qwen3-8B multi-view run (899 of 1367 source groups, $\sim$66\%) -- the corresponding
numbers may shift once these complete. Finally, LoRA fine-tuning was only run on
Qwen3-4B, so it is unknown whether the CnBD-generalization and IPI-collapse pattern holds
at larger scale.

\textbf{Future work.} We plan to use the Robo2VLM-1 pipeline to generate a larger-scale
dataset for more robust cross-difficulty testing; to test a list-position hypothesis for
\texttt{phase\_success} by swapping \texttt{label\_phase} text for the same-position
phase from a different scene (if the answer is unchanged, this would confirm a
position-only cue rather than temporal grounding); to run LoRA on a larger model
(Qwen3-8B) to check whether the gain-decomposition pattern holds at scale; and to run an
error taxonomy on the CoT transcripts to confirm the ``correct description, wrong
conclusion'' failure mode with concrete examples.


{
    \small
    \bibliographystyle{ieeenat_fullname}
    \bibliography{main}
}

% \clearpage
\setcounter{page}{1}
\maketitlesupplementary

\section{Dataset Statistics}
\label{sec:suppl-dataset}

\begin{table}[h]
  \centering
  \caption{Dataset statistics. The LoRA split uses a scene-level partition with no
  scene overlap between train, val, and test.}
  \label{tab:dataset_stats}
  \begin{tabular}{lrrr}
    \toprule
    & \textbf{Train} & \textbf{Val} & \textbf{Test} \\
    \midrule
    Scenes      & 20    & 7   & 20   \\
    Questions   & 2535  & 983 & 2706 \\
    \midrule
    \multicolumn{4}{l}{\small Total: 47 scenes, 6224 questions} \\
    \bottomrule
  \end{tabular}
\end{table}

\section{LoRA Fine-Tuning Details}
\label{sec:suppl-lora}

\begin{table}[h]
  \centering
  \caption{LoRA fine-tuning hyperparameters for Qwen3-4B.}
  \label{tab:lora_hparams}
  \begin{tabular}{ll}
    \toprule
    \textbf{Hyperparameter} & \textbf{Value} \\
    \midrule
    LoRA rank               & 16 \\
    LoRA alpha              & 32 \\
    LoRA dropout            & 0.05 \\
    Target modules          & q, k, v, o projections \\
    Epochs                  & 3 \\
    Learning rate           & $2\times10^{-4}$ \\
    Batch size              & 1 (grad.\ accum.\ 8; effective 8) \\
    Max sequence length     & 1024 tokens \\
    Seed                    & 42 \\
    \midrule
    Train scenes / questions  & 20 / 2535 \\
    Val scenes / questions    & 7 / 983 \\
    Test scenes / questions   & 20 / 2706 \\
    \bottomrule
  \end{tabular}
\end{table}

The train/val/test split is scene-level with no overlap, preventing the model from
memorizing specific visual scenes. We compare against zero-shot Qwen3-4B on the identical
test split, giving an apples-to-apples BASE vs.\ +LoRA comparison.
  % uncomment if supplementary needed

\end{document}
